El código Mermaid devuelto inicialmente por la IA fue sintácticamente válido y representó correctamente la actualización dinámica sin recarga de página y el contador de resultados coincidentes. No obstante, en la revisión técnica se detectaron tres aspectos que se corrigieron manualmente:
\begin{itemize}[leftmargin=*,nosep]
\item \textbf{Endpoint y parámetros de consulta en inglés:} La IA generó \texttt{GET /api/products?category=Vestidos\&size=M}. Se sustituyó por la ruta del glosario común: \texttt{GET /api/productos (categoria, talla)}.
\item \textbf{Sustitución de SQL plano por métodos de acceso a datos y de dominio:} La IA planteó una consulta sobre una única tabla \texttt{products} ignorando las entidades \texttt{CATEGORIA} y \texttt{STOCK\_TALLA}. Se reemplazó por la llamada al método de repositorio \texttt{findProductosByCategoriaAndTalla(categoria, talla)} (que devuelve \texttt{Producto[]} y \texttt{totalCoincidentes}) y el método de dominio \texttt{Producto.estaDisponible()}.
\item \textbf{Alineación con el caso de uso individual (\texttt{cu\_us06}):} Para cubrir la relación \texttt{«extend»} \emph{Limpiar filtros} de \texttt{cu\_us06}, se añadió el bloque opcional \texttt{opt limpiarFiltros()} que invoca \texttt{GET /api/productos} y \texttt{findAllProductosConStock()}, restableciendo el listado completo y el contador, además de configurar \texttt{mirrorActors: false}.
\end{itemize}
